\documentclass[11pt,a4paper]{article}

% --- Encoding and font ---
\usepackage[utf8]{inputenc}
\usepackage[T1]{fontenc}
\usepackage{lmodern}  % Better fonts

% --- Math support ---
\usepackage{amsmath, amssymb}

% --- Page layout ---
\usepackage[margin=2.5cm]{geometry}

% --- Better lists ---
\usepackage{enumitem}

% --- Hyperlinks (optional) ---
\usepackage[colorlinks=true, linkcolor=blue, urlcolor=blue]{hyperref}

% --- Title ---
\title{Tensor Omics planning}
\author{Bobitos}
\date{\today}

\begin{document}

\maketitle

\section{Google Slides for illustrations of figures}

See \href{https://docs.google.com/presentation/d/1T9e2uZamu__GOxcbz7fM81QUgAlvKp5JjUTwnGWbz3w/edit?usp=sharing}{here}, please.

\section{Relative Axes Expression (RAE) Projection}

\subsection{Method: Projection onto the Diagonal-Orthogonal Plane}

Given gene expression vectors in an $n$-dimensional semantic space (e.g., tissues, conditions), we project them onto the hyperplane orthogonal to the space diagonal to analyze relative expression across axes.

Let:
\begin{itemize}
  \item $\mathbf{x} \in \mathbb{R}^n$: the expression vector of a gene (e.g., paralog or ortholog),
  \item $\mathbf{d} = \frac{1}{\sqrt{n}}(1, 1, \dots, 1)$: the unit space diagonal vector.
\end{itemize}

We define the projection of $\mathbf{x}$ onto the plane orthogonal to $\mathbf{d}$ as:
\[
\mathbf{x}_{\perp} = \mathbf{x} - (\mathbf{x} \cdot \mathbf{d}) \mathbf{d}
\]

Let $\mathbf{o}_{\perp}$ and $\mathbf{p}_{\perp}$ be the projections of the ortholog centroid and the paralog outlier, respectively. Then:
\begin{align*}
\Delta \mathbf{v} &= \mathbf{p}_{\perp} - \mathbf{o}_{\perp} \quad \text{(difference vector)} \\
\hat{\Delta \mathbf{v}} &= \frac{\Delta \mathbf{v}}{\|\Delta \mathbf{v}\|} \quad \text{(normalized difference vector)}
\end{align*}

To analyze shifts in direction, we define the \textbf{clock hand vectors}:
\[
\mathbf{v}_{\text{hand}} = \mathbf{x}_{\perp} - \mathbf{d}_{\perp}
\]
which point from the diagonal anchor to the projected expression vector.

The \textbf{angle $\varphi$} between clock hands (e.g., of a paralog and its ortholog) quantifies the shift in \textbf{tissue preference} or semantic emphasis.

\subsection{Biological Interpretability}

\begin{itemize}
  \item \textbf{Angle $\varphi$}: Measures the rotation in tissue (or axis) preference. A change in $\varphi$ indicates that the expression emphasis has shifted between semantic dimensions (e.g., from one tissue to another).
  
  \item \textbf{Difference vector $\Delta \mathbf{v}$}: Represents the raw shift in relative expression across semantic axes.

  \item \textbf{Normalized difference vector $\hat{\Delta \mathbf{v}}$}: Reflects the relative (percental) change in semantic emphasis, independent of expression magnitude.

  \item \textbf{Subfunctionalization}: Can be detected if the clock hand vectors of multiple paralogs sum approximately to the ortholog’s clock hand:
  \[
  \mathbf{v}_{\text{hand}}^{(p_1)} + \mathbf{v}_{\text{hand}}^{(p_2)} \approx \mathbf{v}_{\text{hand}}^{\text{orth}}
  \]

  \item \textbf{Neofunctionalization}: Can be inferred when a paralog expresses in a semantic dimension (e.g., tissue) where the ortholog showed no expression, observable as clock hands extending into previously unoccupied angular sectors.
\end{itemize}

\subsection{Terminology}

\begin{itemize}
  \item \textbf{RAE}: Relative Axes Expression — the name for this projection-based analysis method, generalizable across omics types and semantic spaces.
  \item \textbf{Plane $\Pi_{\perp}$}: The plane orthogonal to the space diagonal $\mathbf{d}$.
  \item \textbf{Clock hand vector}: A vector drawn from the projected diagonal anchor to the projected expression vector.
  \item \textbf{Angle $\varphi$}: Captures the directional change in expression emphasis between genes.
\end{itemize}

\subsection{Clock Plots for Time Trajectory Visualization}

To visualize gene expression trajectories over time in semantic vector space (e.g., tissue or condition axes), we introduce \textbf{clock plots} in the plane $\Pi_{\perp}$ orthogonal to the space diagonal. Each clock hand represents the projected expression vector $\mathbf{x}_{\perp}^{(t)}$ at a specific time point $t$. The twelve o’clock position is defined by a reference vector, typically:
\begin{itemize}
  \item a projected axis (e.g., the projected \textit{leaf} or \textit{seed} axis), or
  \item the direction orthogonal to the conserved mean trajectory (in divergence studies).
\end{itemize}

The clock plot offers intuitive interpretation:
\begin{itemize}
  \item The \textbf{angle between hands} shows directional changes in semantic emphasis over time (e.g., a shift from one tissue to another).
  \item The \textbf{length of the hand} corresponds to the magnitude of deviation from uniform expression across axes.
  \item Time points can be annotated directly on the clock to visualize expression trajectories as semantic rotations.
\end{itemize}

\paragraph{Helper Distance-Time Plot.}
To complement clock plots, we employ a 2D helper plot:
\begin{itemize}
  \item The $x$-axis represents time along the trajectory.
  \item The $y$-axis shows the Euclidean distance between the expression vector at time $t$ and the projected ortholog centroid $\mathbf{o}_{\perp}$.
\end{itemize}

Together, these plots convey both the \textbf{directional trajectory} and the \textbf{divergence magnitude} over time—offering a concise geometric view of neofunctionalization, subfunctionalization, or dynamic tissue shifts in response to experimental conditions.


\paragraph{Projected Semantic Axes as Reference Frame.}
To enhance interpretability of clock plots, the original semantic axes (e.g., tissue, condition, or omics dimensions) are projected into the diagonal-orthogonal plane $\Pi_{\perp}$. These projected axes form a \textbf{warped coordinate star} radiating outward from the origin of the clock.

Let $\mathbf{a}_1, \dots, \mathbf{a}_n$ denote the $n$ standard basis vectors (unit vectors along each semantic axis). We define their projections onto $\Pi_{\perp}$ as:
\[
\mathbf{a}_{i,\perp} = \mathbf{a}_i - (\mathbf{a}_i \cdot \mathbf{d}) \mathbf{d}
\]
These projected axes are visualized as thin rays or arrows in the clock plot, forming a \textbf{star-like reference system}. Due to the off-center perspective created by the orthogonal projection, the angles between the projected axes are generally no longer 90° and become \textbf{visually warped}. This geometric warping encodes the semantic structure of the space and allows users to interpret trajectory movements in relation to the original biological meaning of the axes.

\paragraph{User-Defined Anchor Axis as Twelve O’Clock.}
To orient the plot meaningfully, a user-defined semantic axis is selected as the \textbf{twelve o’clock anchor}. Its projected vector $\mathbf{a}_{\text{anchor},\perp}$ defines the upward direction of the clock. All angle calculations (e.g., $\varphi$ between time points or between paralogs and orthologs) are computed relative to this anchor axis.

This design enables intuitive biological interpretation, e.g.:
\begin{itemize}
  \item setting \textit{leaf} as twelve o’clock lets one interpret early/late developmental tissue shifts as clockwise or counterclockwise movement,
  \item selecting the direction orthogonal to the conserved reference trajectory highlights deviations from evolutionary conservation.
\end{itemize}

\medskip

Together, the warped semantic star and the twelve o’clock anchor provide a geometrically faithful yet biologically meaningful visual frame for interpreting expression dynamics and functional divergence across time, gene families, or omics types.

\end{document}
